\documentclass[a4paper,11pt]{article}
\usepackage[utf8]{inputenc}
\usepackage[T1]{fontenc}
\usepackage[french]{babel}
\usepackage[left=2cm,right=2cm,top=2cm,bottom=2cm]{geometry}
\usepackage{xcolor}
\usepackage{hyperref}
\usepackage{fontawesome5}
\usepackage{parskip}

% --- Couleurs Wavestone ---
\definecolor{primary}{RGB}{35, 55, 123}
\hypersetup{colorlinks=true, urlcolor=primary}

\begin{document}

% --- EXPÉDITEUR ---
\noindent
\begin{minipage}[t]{0.5\textwidth}
    \textbf{Loïc Aron MBASSI EWOLO}\\
    Élève-Ingénieur Bac+5 -- Double diplôme ENSEA\\[3pt]
    \faMapMarker\ Cergy, France\\
    \faPhone\ (+33) 7 59 52 33 98\\
    \faEnvelope\ \href{mailto:loicaron.mbassiewolo@ensea.fr}{loicaron.mbassiewolo@ensea.fr}
\end{minipage}
\hfill
% --- DATE ---
\begin{minipage}[t]{0.4\textwidth}
    \raggedleft
    Cergy, le 9 mars 2026
\end{minipage}

\vspace{1.2cm}

% --- DESTINATAIRE ---
\hfill
\begin{minipage}[t]{0.5\textwidth}
    \raggedleft
    \textbf{Wavestone}\\
    Service Recrutement\\
    Puteaux, France
\end{minipage}

\vspace{1cm}

\noindent\textbf{Objet :} Candidature en alternance -- Data, IA \& IoT (Practice TIME)

\vspace{0.8cm}

Madame, Monsieur,

Orchestrer 4 agents IA via RAG et LangChain pour produire des recommandations agricoles multi-critères, puis concevoir des gants IoT instrumentés traduisant la langue des signes en texte : ces deux projets m'ont confronté à la chaîne complète qui va de la captation de données terrain jusqu'à l'aide à la décision par l'IA. C'est cette logique de bout en bout, appliquée aux enjeux industriels et énergétiques, que je souhaite approfondir au sein de l'équipe Data, IA \& IoT de la practice TIME de Wavestone.

En stage à INRIA Saclay (équipe TRIBE, avril à août 2026), je construis un pipeline Python d'analyse automatisée à grande échelle : NLP sur des corpus de politiques de confidentialité (BERT/Transformers), analyse statique de SDK mobiles et élaboration d'un score de confiance interprétable. Ce travail m'a appris à structurer des données hétérogènes, à évaluer la qualité d'un modèle et à restituer des résultats exploitables par des non-spécialistes, compétences directement transposables aux missions d'acculturation data et de structuration de PoC que propose votre offre.

Côté IoT, mon stage au laboratoire IoT de l'École Polytechnique de Yaoundé (TinyML, TF Lite, Arduino/Raspberry Pi) et le projet Harmony Gloves (STM32, capteurs IMU, fusion capteurs + vision) m'ont donné une compréhension concrète des contraintes embarquées : mémoire limitée, temps réel, intégration hardware/software. La construction d'une plateforme IoT interne pour former les consultants Wavestone est un sujet sur lequel je peux apporter cette expérience terrain. Par ailleurs, le 1\textsuperscript{er} prix au hackathon national CITS (75 équipes), obtenu en optimisant la collecte de déchets urbains par ML sur données IoT ($-$20\% de coûts logistiques), illustre ma capacité à transformer des données capteurs en valeur opérationnelle mesurable.

Je suis disponible dès septembre 2026, à l'issue de mon stage INRIA, pour une alternance de 12 mois. La double culture technique (IA, IoT, data) et conseil que propose Wavestone correspond à l'environnement dans lequel je veux évoluer : des sujets concrets, un accompagnement structuré, et des livrables à impact pour les clients.

Je reste à votre disposition pour un entretien afin de discuter de ma candidature.

Je vous prie d'agréer, Madame, Monsieur, l'expression de mes salutations distinguées.

\vspace{0.8cm}

\textbf{Loïc Aron MBASSI EWOLO}\\
\textit{Élève-Ingénieur ENSEA / Polytechnique Yaoundé}

\end{document}
